\clearpage
\fancyhead[L]{\small BEIMA. DE ANNULO SATURNI.}
\hypertarget{native-659-middle}{}
\begin{center}\small Göttingische gelehrte Anzeigen. No. 194. pp. 1937..1938. 5 December 1842.\\[3pt]\rule{.72\linewidth}{.3pt}\end{center}

\emph{De annulo Saturni commentatus est} \textsc{Elto Martini Beima}, \emph{mus. hist. nat. publ. conservator.} Leyden 1842. Published by S. and J. Luchtmans. 252 pages in quarto, with 4 copperplates.

In this monograph on one of the most interesting objects in our solar system, there has been gathered with great diligence everything concerning the history of the discovery of Saturn's ring, and everything that astronomers have so far observed, calculated or conjectured about its position, shape, dimensions, mass, physical constitution, rotation and origin, as well as about the forces that sustain it and the phenomena it presents to its primary planet. The work can lay no claim to any rank higher than that of a compilation prepared with diligence and drawn from the sources, but one can grant it this praise if one makes allowances for its very faulty and careless language, which in places is scarcely intelligible without consulting the sources from which it was translated.
